\documentclass[11pt]{article}
\usepackage[margin=1in]{geometry}
\usepackage{amsmath,amssymb,amsthm,mathtools}
\usepackage{booktabs,tabularx,longtable}
\usepackage{enumitem}
\usepackage[expansion=false]{microtype}
\usepackage[hidelinks]{hyperref}
\usepackage{cleveref}

% ---------------------------------------------------------------------------
% Status conventions (CLAUDE.md rule 0.1.6).
%   Proposition (proved): a complete proof is in proofs/ and is cited here.
%       It is NOT called a theorem until the proof has been checked line by
%       line by someone other than its author and the check is recorded.
%   Conjecture: no complete proof exists. Each conjecture names the
%       experiment that tests it.
%   Definition: not a claim; carries a quantitative test and a proposed
%       threshold. Every threshold below is PROPOSED. Thresholds are fixed
%       only in prereg/PHASE_X.md before the tag prereg-phase-X and are
%       calibrated on planted-answer systems (Phase 1), never on real data.
% ---------------------------------------------------------------------------

\theoremstyle{definition}
\newtheorem{definition}{Definition}[section]
\newtheorem{assumption}[definition]{Assumption}
\theoremstyle{plain}
\newtheorem{proposition}[definition]{Proposition}
\newtheorem{conjecture}[definition]{Conjecture}
\newtheorem{lemma}[definition]{Lemma}
\theoremstyle{remark}
\newtheorem{remark}[definition]{Remark}

\newcommand{\X}{\mathcal{X}}
\newcommand{\Y}{\mathcal{Y}}
\newcommand{\B}{\mathcal{B}}
\newcommand{\Sites}{\mathcal{S}}
\newcommand{\E}{\mathbb{E}}
\renewcommand{\Pr}{\mathbb{P}}
\newcommand{\pa}{\mathrm{Pa}}
\newcommand{\Hent}{H}
\newcommand{\G}{\mathcal{G}}
\newcommand{\Part}{\pi}
\newcommand{\VI}{\mathrm{VI}}
\newcommand{\IIA}{\mathrm{IIA}}
\newcommand{\status}[1]{\par\noindent\textbf{Status.} #1}
\newcommand{\predict}[1]{\par\noindent\textbf{Measurable prediction.} #1}
\newcommand{\experiment}[1]{\par\noindent\textbf{Experiment.} #1\par\medskip}

\title{The Structure of Understanding: theory core\\[2pt]
\large Sections 2.1--2.4 (Phase 0) with Phase 1 amendments}
\author{WhiteBox E\&E (authors and affiliation to be set by Arjun Paradkar)}
\date{Draft of 2026-10-23. No result in this document is evidence; planted-answer checks are estimator tests only.}

\begin{document}
\sloppy
\maketitle

\begin{abstract}
This document fixes the objects that the rest of the programme measures. It
defines a concept inside a system as a variable of internal state that passes
three tests (quotient, derivation, causal use), defines the concept graph and
the equivalence it is defined up to, defines generality, principles, depth of
understanding and structural memorization, and defines the three moves a new
input can require (lower, sideways, diagonal) together with creation and false
creation. Each definition carries a quantitative test and a proposed threshold.
Each claim carries its status. Eight propositions have complete proofs in
\texttt{proofs/} that have not yet been independently checked; every other
claim is a conjecture with a named experiment. Nothing here is evidence; the
evidence is the pre-registered programme in Sections 5--6 of the master
prompt.
\end{abstract}

\tableofcontents

%============================================================================
\section{Setting (master prompt \S2.1)}
%============================================================================

\begin{definition}[System, sites and states]\label{def:system}
A \emph{system} is a measurable map $f\colon\X\to\Y$ computed through a finite
set of \emph{sites} $\Sites$ partially ordered by a computation order
$\prec$ (for a network: (layer, position, component) triples ordered by the
forward pass; for a brain: (region, time window) pairs ordered by time).
Each site $s$ carries a state $h_s(x)\in\mathbb{R}^{d_s}$. We fix once and for
all a linear extension of $\prec$, written $s_1<s_2<\dots<s_{|\Sites|}$, and
write $s<t$ for it. The input $x$ is treated as the state of a source site
$s_0$, and the output as the state of a sink site $s_\star$.
\end{definition}

\begin{definition}[Task distribution and behaviours]\label{def:behaviour}
A task distribution $P$ on $\X$ is given. A \emph{behaviour} is a pair
$b=(P_b,r_b)$ of an input distribution $P_b\ll P$ and a readout
$r_b\colon\Y\to\mathcal{R}_b$ into a finite set or a metric space (for a
language model: the answer token of a task family; for an animal: the choice
on a trial type). A \emph{behaviour suite} is a finite family
$\B=\{b_1,\dots,b_m\}$ with weights $w_b\ge0$, $\sum_b w_b=1$. For each $b$ we
fix a divergence $d_b$ on the outputs relevant to $r_b$ (total variation
between output distributions, or $0/1$ disagreement of $r_b$).
\end{definition}

\begin{definition}[Interventions and access levels]\label{def:intervention}
For a site $s$, a \emph{realiser} is a map $R\colon\mathbb{R}^{d_s}\to
\mathbb{R}^{k}$ together with a right inverse that writes $R$-coordinates back
into the state while leaving the complementary part fixed (a linear subspace
with its orthogonal projection, as in distributed alignment search; or a set
of sparse-dictionary features whose activations are swapped while the
reconstruction error is kept). For inputs $x,x'$ the \emph{realiser patch}
$f_{R\leftarrow x'}(x)$ is the output of the run on $x$ in which the
$R$-coordinates of $h_s(x)$ are replaced by those of $h_s(x')$ and all later
computation is recomputed. Patches at several realisers are applied jointly.
A result obtained with patches is at access level (C). At level (B) a patch is
replaced by its first-order Jacobian estimate; at level (D) by the strongest
perturbation the data allow (stimulation, optogenetics, or, failing that,
decoding plus a natural-experiment contrast); at level (A) by counterfactual
input edits. Every statement below is written for level (C); its weaker
versions inherit the level of the instrument and must say so.
\end{definition}

\begin{remark}[Nulls attached to the setting]
Every quantity defined below is computed on three systems in addition to $f$:
an untrained copy $f^{0}$ with the same architecture (master prompt \S1.4),
a feature-shuffle null that keeps marginals, and random realisers of matched
dimension and norm. A quantity earns a claim only through its contrast with
these nulls.
\end{remark}

%============================================================================
\section{Concept (master prompt \S2.2)}
%============================================================================

A candidate concept is a triple $(s,R,\varphi)$ of a site, a realiser at that
site, and a \emph{value map} $\varphi\colon\mathbb{R}^k\to V$ into a finite set
or a metric space. Its value on an input is
\[
C(x)=\varphi\bigl(R(h_s(x))\bigr).
\]
Candidates are proposed by sparse dictionaries, by the null space of the
pullback metric $g_s=J_s^\top J_s$ of the downstream map (master prompt
\S1.9), and by probes for known variables. Proposal confers nothing; only the
three conditions below do.

\begin{definition}[Partition of a concept]
The \emph{partition} of $C$ is $\Part_C=\{C^{-1}(v):v\in V\}$, taken modulo
$P$-null sets. For continuous $V$ with metric $\rho$, fibres are replaced by
$\delta$-fibres $\{(x,x'):\rho(C(x),C(x'))\le\delta\}$ at a resolution
$\delta$ fixed in pre-registration.
\end{definition}

\subsection{Condition 1: quotient}

\begin{definition}[Quotient defect and compression]\label{def:quotient}
For a behaviour $b$, the \emph{within-fibre swap effect} and the
\emph{unconditioned swap effect} are
\[
A_b^{=}(C)=\E_{x,x'\sim P_b\,:\,C(x)=C(x')}\,
  d_b\bigl(f(x),f_{R\leftarrow x'}(x)\bigr),\qquad
A_b(C)=\E_{x,x'\sim P_b}\,
  d_b\bigl(f(x),f_{R\leftarrow x'}(x)\bigr),
\]
and the \emph{quotient defect} is
$\delta_Q(C)=\sum_b w_b\,A_b^{=}(C)\big/\sum_b w_b\,A_b(C)$ over the behaviours
with $A_b(C)>\eta$. The \emph{compression ratio} is
$\kappa(C)=\Hent(C)/\Hent(X)$, with $\Hent(X)=\log N_{\mathrm{eff}}$ for an
empirical input set of effective size $N_{\mathrm{eff}}$.
\end{definition}

\noindent\textbf{Quantitative test.} $C$ passes the quotient condition when
(i) the realiser matters, $\sum_b w_bA_b(C)\ge\eta$ with $\eta$ the
95th percentile of the same quantity for random realisers of matched
dimension and norm;
(ii) $\delta_Q(C)\le\varepsilon_Q$ with the upper end of its 95\% bootstrap CI
at most $2\varepsilon_Q$;
(iii) $\delta_Q(C)$ is below the 5th percentile of $\delta_Q$ for random
partitions of $X$ with the same block-size profile as $\Part_C$
(permutation test, Holm-corrected across candidates);
(iv) $\kappa(C)\le\kappa_{\max}$ and $\Hent(C)\ge h_{\min}$.
\textbf{Proposed thresholds:} $\varepsilon_Q=0.10$, $\kappa_{\max}=0.25$,
$h_{\min}=0.5$ bit.

\begin{remark}
Condition 1 says that, inside its realiser, nothing but the value of $C$ is
visible to what follows. It does not say that the downstream computation
depends on $h_s$ only through $C$: other variables at the same site may be
used too, and they are other candidates. Swaps that change the value of $C$
are the subject of Condition 3. Chalupka, Perona and Eberhardt's causal
coarsening of an input space is the closest prior construct; the difference is
that the partition here is of internal states and is tested through realiser
swaps.
\end{remark}

\subsection{Condition 2: derivation}

Let $\mathcal{K}_{<s}$ be the set of concepts already accepted at sites
strictly before $s$, together with the input variables (tokens, pixels,
stimulus labels), which are the \emph{input nodes}.

\begin{definition}[Sufficient and minimal parent sets]\label{def:parents}
A set $S\subseteq\mathcal{K}_{<s}$ is \emph{$\varepsilon$-sufficient} for $C$
if $\Hent(C\mid S)\le\varepsilon\,\Hent(C)$, estimated as the held-out
conditional cross-entropy of a sparse predictor $\hat\psi$ of $C$ from $S$.
It is \emph{minimal} if no proper subset is $2\varepsilon$-sufficient. The
\emph{parent set} $\pa(C)$ is the minimal $\varepsilon_D$-sufficient set
returned by sparse conditional prediction, and the \emph{derivation map} is
$\psi_C=\hat\psi$, so that $C\approx\psi_C(\pa(C))$.
\end{definition}

\noindent\textbf{Quantitative test.} $C$ passes the derivation condition when
(i) a minimal $\varepsilon_D$-sufficient set exists with $|\pa(C)|\le k_{\max}$;
(ii) the parent set is stable: mean pairwise Jaccard index at least $J_{\min}$
across at least five data resamples and, where they exist, five seeds;
(iii) each edge is interventionally confirmed: patching the realiser of a
parent $C_i$ from $x'$ into $x$ moves $C$ to $\psi_C$ evaluated at the patched
parent values on at least a fraction $\tau_E$ of pairs, and at least $0.2$
above random realisers.
\textbf{Proposed thresholds:} $\varepsilon_D=0.05$, $k_{\max}=8$,
$J_{\min}=0.8$, $\tau_E=0.9$.

\begin{proposition}[Derivation is refinement in the partition lattice;
proved, \texttt{proofs/P1\_partition\_lattice.tex}]\label{prop:lattice}
Order partitions of $X$ by refinement, $\Part\le\Part'$ if every block of
$\Part$ lies inside a block of $\Part'$, and write $\bigwedge$ for the common
refinement (meet). In the exact case $\varepsilon_D=0$ and for finite $X$,
$C=\psi(C_1,\dots,C_k)$ for some map $\psi$ if and only if
$\bigwedge_{i}\Part_{C_i}\le\Part_C$. Consequently $S$ is an exact parent set
of $C$ if and only if $S$ is an inclusion-minimal subset of
$\mathcal{K}_{<s}$ whose meet refines $\Part_C$.
\end{proposition}
\predict{None by itself; it fixes what the estimator must recover and is the
basis of the planted-answer check.}
\experiment{Phase 1, planted system V1 (hand-built digit, column-sum, carry,
result network): the estimated parent sets must equal the planted ones on
every node at the pre-registered recovery threshold.}

\begin{proposition}[Uniqueness of the minimal parent set under positivity;
proved, \texttt{proofs/P2\_parent\_uniqueness.tex}]\label{prop:unique}
Let $U=\mathcal{K}_{<s}$ be finite-valued and suppose the joint law of $U$
under $P$ gives positive probability to every combination of values in the
product of the observed value sets of its members (positivity of the
candidates; $C$ itself may be a deterministic function of $U$). Then the
inclusion-minimal set $S\subseteq U$ with $C\perp U\setminus S\mid S$ is
unique.
\end{proposition}
\predict{Where positivity holds, parent sets estimated on disjoint data halves
coincide up to sampling error; where it fails, they may differ, and the
difference is located on value combinations that never co-occur.}
\experiment{Phase 1, planted systems V1 and V3; on V3 (components seen
separately, never together) parent-set stability must drop exactly at the
withheld combinations.}

\begin{remark}[Positivity is a real restriction]
Positivity fails whenever the candidate set contains a node together with a
function of it, for example a node and its own parents, so it cannot be
assumed for all of $\mathcal{K}_{<s}$. With $C=C_3=C_1\oplus C_2$ both
$\{C_3\}$ and $\{C_1,C_2\}$ are minimal sufficient sets. Until a canonical
candidate set is fixed (one option: the latest-site antichain that is
sufficient), the programme reports all minimal parent sets and every
downstream quantity as a range over the resulting graphs. Fixing the rule is
a Phase 1 design decision, tested on V1 before any real model.
\end{remark}

\begin{proposition}[Derivation maps are unidentified on unseen combinations
and identified by interchange; proved,
\texttt{proofs/P3\_unseen\_combinations.tex}]\label{prop:unseen}
(a) If a combination $t$ of parent values has probability zero under $P$,
then for any $\psi$ that agrees with $C$ $P$-almost surely there is a $\psi'$
that also agrees with $C$ $P$-almost surely and differs from $\psi$ at $t$.
(b) Suppose the parent realisers are \emph{modular}: there exist inputs
$x,x^{(1)},\dots,x^{(k)}$ such that the joint patch of the realiser of $C_i$
from $x^{(i)}$ into the run on $x$ produces parent values $t$ and leaves the
computation from the parents to $C$ unchanged in form. Then the value the
system assigns to $C$ at $t$ is observed directly, so the system's derivation
map is identified at $t$ by the interchange intervention.
\end{proposition}
\predict{Observational fits cannot predict sideways behaviour; interchange
tests on the new combination can. This is the formal content of the sideways
prediction in \S\ref{sec:moves}.}
\experiment{Phase 2--3, K3: on sideways inputs, the agreement between the
interchange-identified $\psi$ and the correct answer predicts correctness
better than any observational fit of $\psi$.}

\subsection{Condition 3: causal use}

\begin{definition}[Interchange intervention accuracy]\label{def:iia}
Let $D_C$ be the \emph{downstream derivation} of $C$: the sub-graph of the
current concept graph from $C$ to the output, with its derivation maps, which
defines a high-level model $F_{\G}$ in the sense of causal abstraction
(Geiger, Lu, Icard and Potts). For a pair $(x,x')$ with $C(x)\neq C(x')$ the
predicted counterfactual output is $\hat y(x;C\leftarrow C(x'))=F_{\G}(x;
\mathrm{do}(C=C(x')))$. The interchange intervention accuracy is
\[
\IIA(C)=\sum_b w_b\,\Pr_{x,x'\sim P_b\,:\,C(x)\ne C(x')}
\Bigl[r_b\bigl(f_{R\leftarrow x'}(x)\bigr)=r_b\bigl(\hat y(x;C\leftarrow C(x'))\bigr)\Bigr].
\]
A graph-free variant replaces $\hat y$ by the system's output on a matched
counterfactual input $x^\ast$ that agrees with $x$ on every other parent of the
children of $C$ and has $C(x^\ast)=C(x')$, where such an input exists in the
sample. This variant breaks the circularity between Condition 3 and the graph
on the first pass of discovery.
\end{definition}

\noindent\textbf{Quantitative test.} $C$ passes the causal-use condition when
$\IIA(C)\ge\tau_{\IIA}$, the lower end of its 95\% CI is above
$\IIA$ of random realisers of matched dimension by at least $\Delta_{\IIA}$,
and the same realiser in the untrained copy $f^0$ does not pass. A candidate
whose probe decoding accuracy is at least $0.9$ but that fails this test is
recorded as \emph{represented, not used}.
\textbf{Proposed thresholds:} $\tau_{\IIA}=0.90$, $\Delta_{\IIA}=0.20$.

\begin{definition}[Concept]\label{def:concept}
A \emph{concept} of $f$ (at access level (C)) is a candidate that passes
Conditions 1, 2 and 3. Input nodes are concepts by convention and are exempt
from Conditions 1 and 2.
\end{definition}

\subsection{The concept graph and its equivalence}

\begin{definition}[Concept graph]\label{def:graph}
Run discovery site by site in the order $<$. The concept graph $\G(f)$ has as
nodes the concepts, taken modulo the equivalence below, and an edge
$C_i\to C$ for each $C_i\in\pa(C)$. A node's site is the earliest site at
which a member of its class is accepted.
\end{definition}

\begin{definition}[Equivalence]\label{def:equiv}
Two concepts $C,C'$ are \emph{equivalent} if $\Part_C=\Part_{C'}$ modulo
$P$-null sets; approximately, if the normalised variation of information
$\VI(\Part_C,\Part_{C'})/\log N_{\mathrm{eff}}\le\varepsilon_{\mathrm{merge}}$.
Two graphs are equivalent if there is a bijection of nodes that maps each node
to an equivalent node and preserves edges. Equivalence absorbs (i) any
invertible reparametrisation $u\colon V\to V'$ of a value map, since
$\Part_{u\circ C}=\Part_C$, and (ii) the merging of redundant nodes that carry
the same partition at different sites or in different realisers.
\textbf{Proposed threshold:} $\varepsilon_{\mathrm{merge}}=0.05$.
\end{definition}

\begin{proposition}[Invariance under reparametrisation; proved,
\texttt{proofs/P4\_reparametrisation.tex}]\label{prop:reparam}
Every test statistic in Conditions 1--3 ($\delta_Q$, $\kappa$, $\Hent(C)$,
$\Hent(C\mid S)$, $\IIA$) and every quantity in \S\ref{sec:gen}
(generality, grounding, memorization) is unchanged when a value map $\varphi$
is replaced by $u\circ\varphi$ for an invertible $u$, provided the derivation
maps of $C$ and of its children are replaced by $u\circ\psi_C$ and
$\psi_{C''}\circ(\dots,u^{-1},\dots)$ respectively.
\end{proposition}

\begin{proposition}[Well-definedness in the exact case; proved,
\texttt{proofs/P2\_parent\_uniqueness.tex}, Corollary]\label{prop:welldef}
Assume (a) the exact case $\varepsilon_Q=\varepsilon_D=0$ and
$\varepsilon_{\mathrm{merge}}=0$, finite $X$ and finite value sets;
(b) positivity of the candidate parents at every site
(Proposition~\ref{prop:unique}); (c) the set of accepted candidates is fixed. Then
$\G(f)$ is a directed acyclic graph that is unique up to the equivalence of
Definition~\ref{def:equiv}, and it does not depend on the choice of value maps within
each equivalence class.
\end{proposition}

\begin{conjecture}[Stability of the approximate graph]\label{conj:stable}
Suppose the exact-case assumptions hold for some limiting system and that
every accepted concept clears each threshold with margin $\mu>0$ and every
rejected candidate fails some threshold by margin $\mu$. Then for
$\varepsilon_Q,\varepsilon_D,\varepsilon_{\mathrm{merge}}$ small relative to
$\mu$ and for sample size $n\to\infty$, the estimated graph equals the
exact-case graph up to equivalence with probability tending to one.
\end{conjecture}
\status{Conjecture. Obstruction to a proof: the sparse-prediction estimator
of $\pa(C)$ needs a margin condition analogous to the beta-min condition of
sparse regression, and its form for conditional entropy is not yet stated.}
\predict{Graphs estimated on disjoint data halves, and on seeds of the same
training recipe, agree up to equivalence at a rate that rises with $n$.}
\experiment{Phase 1 (V1--V4), recovery vs sample size; Phase 2 (K1),
split-half graph agreement on Llama, Qwen, Gemma and OLMo.}

\begin{remark}[Dependence on the linear extension]
Acyclicity holds by construction because edges respect the fixed linear
extension $<$. Whether the graph depends on the choice of extension when sites
are incomparable under $\prec$ (parallel attention heads, parallel brain
regions) is open; it is recorded as Conjecture~\ref{conj:stable}'s companion
and tested by recomputing the graph under two extensions.
\end{remark}

\subsection{Relation to Formal Concept Analysis}

A formal context is a triple $(G,M,I)$ of objects, attributes and an
incidence relation; with $A'=\{m:\forall g\in A,\ gIm\}$ and
$B'=\{g:\forall m\in B,\ gIm\}$, a formal concept is a pair $(A,B)$ with
$A'=B$ and $B'=A$ (Wille 1982; Ganter and Wille 1999). We attach to a site the
\emph{downstream context} $\mathbb{K}_s=(X,M_s,I)$, whose attributes $M_s$ are
binary downstream readouts $m\colon X\to\{0,1\}$ (thresholded behaviour
readouts and thresholded later-site features), with $xIm$ iff $m(x)=1$. Say
that $C$ is an \emph{exact quotient for $M_s$} if every $m\in M_s$ is constant
on every block of $\Part_C$.

\begin{proposition}[Quotient condition and formal concepts; proved,
\texttt{proofs/P5\_fca.tex}]\label{prop:fca}
Let $C$ be an exact quotient for $M_s$.
(a) Every extent of $\mathbb{K}_s$ is a union of blocks of $\Part_C$.
(b) For every $x$, the block $[x]_C$ is contained in the object-concept extent
$\{x\}''$.
(c) If $M_s$ is closed under complement ($m\in M_s\Rightarrow 1-m\in M_s$,
the nominal scaling of FCA), then the object-concept extents
$\{\{x\}'':x\in X\}$ are exactly the blocks of the kernel partition
$\Part_{M_s}$ of $x\mapsto(m(x))_{m\in M_s}$; $\Part_C\le\Part_{M_s}$; and $C$
is the coarsest exact quotient for $M_s$ if and only if its blocks are the
object-concept extents.
\end{proposition}

\begin{conjecture}[Approximate version]\label{conj:fca}
If $C$ has quotient defect $\delta_Q(C)\le\varepsilon$ with respect to a
behaviour suite whose thresholded readouts form $M_s$, then every extent of
$\mathbb{K}_s$ differs from a union of $\Part_C$-blocks on a set of
$P$-probability at most $c\,\varepsilon$, with $c$ depending only on the
thresholding margin of the readouts.
\end{conjecture}
\status{Conjecture. Obstruction: $\delta_Q$ is a ratio of averaged
divergences, so it bounds the mass of disagreeing pairs only after the
readouts are thresholded with a margin.}
\predict{On a model with a known graph, FCA extents of the downstream context
are, to within $c\,\varepsilon$, unions of the fibres of the accepted
concepts; on the untrained copy they are not.}
\experiment{Phase 1, V1 and V2; reported as a check on Condition 1, never as
evidence for a claim.}

\begin{remark}[What the FCA relation does and does not say]
Proposition~\ref{prop:fca} shows that the quotient condition is a one-site, interventional
form of FCA's closure: concepts that pass Condition 1 induce saturated extents,
and under nominal scaling the coarsest such concept is FCA's object-concept
partition. FCA by itself has neither a derivation relation (its lattice order
is subsumption of extents, not computation of one variable from others) nor a
causal-use test, so Conditions 2 and 3 are not consequences of it.
\end{remark}

%============================================================================
\section{Generality, principles and depth (master prompt \S2.3)}
\label{sec:gen}
%============================================================================

\begin{definition}[Active derivation of an input]\label{def:active}
For an input $x$ in the support of $P_b$, node $C$ is \emph{active on $x$}
if its realiser swap moves the readout on a majority of partners:
\[
a_C(x)=\mathbf{1}\Bigl[\Pr_{x'\sim P_b:\,C(x')\neq C(x)}
\bigl(r_b(f_{R\leftarrow x'}(x))\neq r_b(f(x))\bigr)\ge\alpha\Bigr].
\]
The active derivation $D(x)$ is the sub-graph of $\G(f)$ induced by the
active nodes that are ancestors of the output node, together with the input
nodes. \textbf{Proposed threshold:} $\alpha=0.5$.
\end{definition}

\begin{definition}[Generality]\label{def:generality}
The \emph{weighted generality} of a node is the behaviour-weighted mass of
inputs on whose derivation it is active,
\[
g(C)=\sum_{b\in\B}w_b\,\E_{x\sim P_b}\,a_C(x)\in[0,1],
\]
and the \emph{count generality} is
$g_\#(C)=\#\{b:\E_{P_b}a_C\ge\alpha\}$. The \emph{null-corrected}
generality is $g(C)-q_{95}$, where $q_{95}$ is the 95th percentile of $g$ over
the nodes of the untrained-copy graph (or, if that graph is empty, over random
realisers of matched dimension).
\end{definition}

\begin{definition}[Principle]\label{def:principle}
A \emph{principle at level $\tau$} is a node that is not an input node and has
$g(C)\ge\tau$ and $g(C)>q_{95}$. Write $\Pi_\tau$ for the set of principles.
\textbf{Proposed thresholds:} $\tau\in\{0.25,0.5\}$ pre-registered as two
levels; all results are also reported as curves in $\tau$.
\end{definition}

\begin{definition}[Depth of understanding]\label{def:depth}
Let $\mathcal{P}(x)$ be the set of directed paths in $D(x)$ from an input
node to the output.
\begin{enumerate}[label=(\roman*)]
\item \emph{Thresholded depth.} The behaviour is \emph{grounded at $\tau$ on
$x$} if $\Pi_\tau\cap D(x)$ is a vertex cut separating the input nodes from
the output in $D(x)$.
\item \emph{Structural grounding level.}
$\lambda(x)=\min_{p\in\mathcal{P}(x)}\max_{C\in p}g(C)$, the bottleneck value
of the most weakly grounded path. Then $x$ is grounded at $\tau$ if and only
if $\lambda(x)\ge\tau$ (Proposition~\ref{prop:bottleneck}).
\item \emph{Interventional grounding share.} With $\Pi_\tau$ patched from $x$
into the run on $x'$,
\[
\gamma_\tau(x)=1-\frac{\E_{x'\sim P_b}\,d_b\bigl(f(x),
f_{\Pi_\tau\leftarrow x}(x')\bigr)}{\E_{x'\sim P_b}\,d_b\bigl(f(x),f(x')\bigr)},
\]
the share of the input's effect on the output that passes through principles
($\gamma=1$: the output on $x$ is fixed once the principles are fixed;
$\gamma=0$: principles carry nothing).
\item \emph{Depth of a behaviour.}
$\Lambda_\tau(b)=\Pr_{x\sim P_b}(\lambda(x)\ge\tau)$ (thresholded),
$\bar\lambda(b)=\E_{P_b}\lambda(x)$ and
$\bar\gamma_\tau(b)=\E_{P_b}\gamma_\tau(x)$ (continuous), and the
$\tau$-integrated score $\int_0^1\bar\gamma_\tau(b)\,d\tau$.
\end{enumerate}
\end{definition}

\begin{proposition}[Bottleneck characterisation; proved,
\texttt{proofs/P6\_bottleneck.tex}]\label{prop:bottleneck}
In any finite DAG with node weights $g$, the nodes with $g\ge\tau$ form a
vertex cut between the input nodes and the output if and only if
$\min_{p}\max_{C\in p}g(C)\ge\tau$. Hence (i) and (ii) of
Definition~\ref{def:depth} agree, and $\lambda(x)$ is computable by a bottleneck path
search in time linear in the number of edges after sorting.
\end{proposition}

\begin{conjecture}[Structural grounding implies interventional grounding]
\label{conj:gamma}
Assume $\G(f)$ is \emph{complete} on $x$ (every computational path from the
input to the output passes through a node of $D(x)$), the quotient and
interchange conditions hold exactly, and realisers are modular. If
$\lambda(x)\ge\tau$ then $\gamma_\tau(x)=1$. With approximations,
$1-\gamma_\tau(x)$ is bounded by a constant times
$\varepsilon_Q+(1-\tau_{\IIA})$ plus the completeness defect.
\end{conjecture}
\status{Conjecture. The exact case has a proof sketch (patching the cut fixes
every downstream value by the quotient and interchange conditions); the
sketch is not complete because completeness of $\G(f)$ is not itself
measurable without a residual term, which has not been defined.}
\predict{$\bar\lambda(b)$ and $\bar\gamma_\tau(b)$ are strongly rank-correlated
across behaviours, and their disagreement is largest where the graph's
mediation share is lowest.}
\experiment{Phase 1 planted systems; Phase 2 K1 (arithmetic, where the
completeness defect can be measured against the known algorithm).}

\begin{definition}[Structural memorization]\label{def:memo}
On input $x$, the behaviour is \emph{memorized at $\tau$} if no principle is
active on its derivation, $\mu(x)=\max_{C\in D(x)\setminus\text{inputs}}g(C)<\tau$
(with $\mu(x)=0$ if $D(x)$ has no non-input node). The behaviour $b$ is
memorized if $\Pr_{P_b}(\mu(x)<\tau)\ge1-\varepsilon_M$, and \emph{grounded} if
$\Lambda_\tau(b)\ge1-\varepsilon_M$; otherwise it is \emph{partly grounded}.
The verdict uses only training-distribution inputs and interventions; it uses
no held-out data. \textbf{Proposed threshold:} $\varepsilon_M=0.10$.
\end{definition}

\begin{proposition}[Order of the scores; proved,
\texttt{proofs/P6\_bottleneck.tex}]\label{prop:order}
$\lambda(x)\le\mu(x)$ for every $x$. Hence an input that is grounded at
$\tau>0$ is not memorized at $\tau$, and, provided $\varepsilon_M<1/2$, the
three verdicts (grounded, partly grounded, memorized) are mutually exclusive
for a behaviour.
\end{proposition}

\begin{conjecture}[Depth predicts generalization; the K2 kill test]
\label{conj:k2}
Across held-out inputs of a behaviour, per-input error is a decreasing
function of $\lambda(x)$ and of $\gamma_\tau(x)$, and these scores predict
per-input error with higher AUROC than Bernardi et al.'s cross-condition
generalization score, max-softmax, entropy, margin and sharpness, at the
pre-registered margin.
\end{conjecture}
\status{Conjecture. This is the programme's kill test; if it fails the
programme stops (master prompt \S6.2, K2).}
\predict{AUROC of $\lambda$ for per-input error exceeds every baseline in
\S6.4 with non-overlapping 95\% CIs, on held-out arithmetic, syntactic
agreement and a diagnosis model.}
\experiment{Phase 2, K2.}

\begin{conjecture}[Memorization verdict predicts held-out failure]
\label{conj:memo}
A behaviour judged memorized by Definition~\ref{def:memo} has held-out accuracy no
better than the best input-free baseline, and a behaviour judged grounded has
held-out accuracy within the training-accuracy CI, provided the specimen gate
(training accuracy at ceiling) is met.
\end{conjecture}
\predict{The structural verdict, computed without held-out data, separates
leakage-trained from clean models and planted-memorized from planted-rule
subsets.}
\experiment{Phase 1 V2 (estimator test only); Phase 4 K4 and K9 (claims).}

\begin{remark}[Relation to Bernardi et al.]
Proposition 2.7.1 of the master prompt (cross-condition generalization as the
depth-one case of depth) is attempted in Phase 1, not here. The expected
form is a statement about a single-node graph with a linear value map and
observational data, with the causal-use term only able to lower the score.
\end{remark}

%============================================================================
\section{The three moves and creation (master prompt \S2.4)}
\label{sec:moves}
%============================================================================

\begin{definition}[Training range]\label{def:range}
For a node $C$ with finite $V$, its \emph{value range} $\mathrm{Rg}(C)$ is the
set of values taken on at least $n_{\min}$ training inputs. For continuous
$V$, $\mathrm{Rg}(C)$ is the set of values whose $k$-nearest-neighbour
distance to training values is below the 99th percentile of the same distance
among training values. The \emph{combination range} of $C$ is
$\mathrm{Rg}(\pa(C))$, the set of parent-value tuples seen on at least
$n_{\min}$ training inputs. A realiser state \emph{has a value} if it lies
within the 99th-percentile distance of the realiser states of some fibre.
\textbf{Proposed thresholds:} $n_{\min}=5$, $k=10$.
\end{definition}

\begin{definition}[Lower, sideways and diagonal]\label{def:moves}
Let $x^\ast$ be an input not seen in training, with active derivation
$D(x^\ast)$.
\begin{enumerate}[label=(\roman*)]
\item $x^\ast$ is \emph{lower} if, for every node $C\in D(x^\ast)$, the
realiser state has a value, $C(x^\ast)\in\mathrm{Rg}(C)$, and the parent tuple
$\pa(C)(x^\ast)\in\mathrm{Rg}(\pa(C))$.
\item $x^\ast$ is \emph{sideways} if every node value is in range but at
least one parent tuple is not.
\item $x^\ast$ is \emph{diagonal} if the derivation needs a concept that is
not in $\G(f)$. Operationally: some realiser state on $D(x^\ast)$ has no value,
or the graph's mediation share
$\gamma_{\G}(x^\ast)$ (Definition~\ref{def:depth}(iii) with all non-input nodes patched)
is below $1-\varepsilon_{\mathrm{cov}}$, so part of the input's effect on the
output passes through no node. In benchmarks, diagonal inputs are also
labelled by construction (the task requires a variable that is not a function
of any node).
\end{enumerate}
\textbf{Proposed threshold:} $\varepsilon_{\mathrm{cov}}=0.2$. The
operational classifier is $r.\mathrm{move}(x)$; its agreement with
construction labels is tested at a pre-registered AUROC.
\end{definition}

\begin{conjecture}[Lower inputs are answered correctly]\label{conj:lower}
If $x^\ast$ is lower and the behaviour is grounded, the system's error on
$x^\ast$ is at most its training error plus a term that vanishes as
$\varepsilon_Q,\varepsilon_D,1-\tau_{\IIA}\to0$.
\end{conjecture}
\status{Conjecture. In the exact case it reduces to the statement that every
node's value on $x^\ast$ is computed by maps that were exercised on the same
arguments in training; the obstruction is that ``exercised in training'' does
not imply ``correct'' unless training error is zero on those arguments.}

\begin{conjecture}[Sideways inputs: correct iff interchange passes]
\label{conj:sideways}
If $x^\ast$ is sideways, the system answers correctly if and only if, for
each unseen parent tuple $t$ on $D(x^\ast)$, the interchange intervention that
assembles $t$ from training inputs (Proposition~\ref{prop:unseen}(b)) yields the value
of $C$ prescribed by the task's ground-truth derivation.
\end{conjecture}

\begin{conjecture}[Diagonal inputs fail without creation]\label{conj:diag}
If $x^\ast$ is diagonal, the system fails on it unless a creation event
(Definition~\ref{def:creation}) occurs under exposure to $x^\ast$'s problem.
\end{conjecture}
\predict{For Conjectures \ref{conj:lower}--\ref{conj:diag} jointly: accuracy
orders lower $>$ sideways $>$ diagonal; within sideways, the interchange
verdict predicts correctness at a pre-registered AUROC; within diagonal,
success occurs only with a valid new node.}
\experiment{Phase 3, K3 (moves); Phase 7, K13 (diagonal with creation).
Estimator check on planted system V3 in Phase 1.}

\begin{definition}[Creation]\label{def:creation}
Let $e_1,e_2,\dots$ be an exposure sequence to a new problem (in-context
examples for a language model, training steps for a network, sessions for an
animal), measured in units $u$ (tokens, gradient steps, sessions), and let
$\G_t$ be the concept graph after exposure $t$, recomputed with the same
estimator and thresholds. A \emph{valid new node} at $t$ is a node $C^\ast$ of
$\G_t$ such that
\begin{enumerate}[label=(\roman*)]
\item $C^\ast$ passes Conditions 1--3 in the exposed system;
\item $C^\ast$ is not equivalent to any node of $\G_0$
($\VI$ above $\varepsilon_{\mathrm{merge}}$ to every node);
\item its parents are nodes of $\G_0$, or new nodes that are themselves valid,
so that it is derived from existing nodes;
\item it does not appear, at the same exposure count, under exposure to a
matched control problem that does not require it.
\end{enumerate}
The \emph{creation rate} is the number of valid new nodes per unit $u$; the
\emph{creation reach} of $C^\ast$ is its generality $g_t(C^\ast)$ over the
original suite $\B$ together with the new behaviour.
\end{definition}

\begin{definition}[False creation]\label{def:falsecreation}
A node that appears under exposure is a \emph{false creation} if
(a) it is decodable but not used (passes probe decoding, fails Condition 3);
(b) it is used but not derived (passes Condition 3, fails Condition 2 with
respect to $\G_0$), which makes it a new leaf, i.e.\ memorization of the new
problem;
(c) it is a relabelling (equivalent to a node of $\G_0$, including an
invertible reparametrisation of one); or
(d) it is exposure-generic (it appears equally under the control problem).
Every reported creation count lists the candidates rejected under (a)--(d).
\end{definition}

\begin{conjecture}[Creation predicts success on a diagonal problem]
\label{conj:creation}
Across systems and problems, success on a diagonal problem after exposure is
predicted by the appearance of at least one valid new node whose parents lie
on the derivation of the problem, and the creation rate is positively
associated with the speed of success; false creations do not predict success.
\end{conjecture}
\status{Conjecture; the creation tests run last because they are the most
likely to fail.}
\predict{Logistic regression of success on (valid-node indicator) has a
positive coefficient with CI excluding zero across the K13 problem set; the
same regression on false-creation counts does not.}
\experiment{Phase 7, K13; estimator check on planted system V5 in Phase 1.}


%============================================================================
\section{Phase 1 amendments M1--M9 (estimator design decisions, 2026-10-23)}
\label{sec:amend}
%============================================================================

Planted-answer development (V1--V5, development seeds only) exposed places
where the Phase 0 definitions either fail on a planted answer or leave an
estimator choice open. Each amendment below was fixed before the frozen
validation seeds were run (\texttt{prereg/PHASE\_1.md}) and is part of the
estimator definition from then on. None was chosen on real data.

\begin{enumerate}[label=(M\arabic*)]
\item \textbf{Interventional direct-effect parents.} With parents chosen by
observational minimal sufficiency (Definition~\ref{def:parents}), a planted
rule-plus-exceptions network (V2) gets the output's parent set $\{a,b\}$ (an
input lookup), so its rule node never lies on any derivation and every input
is scored memorized. The parent set is therefore the set of pool nodes $P$
at earlier sites whose \emph{direct} effect on $C$ (swap $P$ from partners
with a different $P$ value, every pool node at a site strictly between held
at its clean value) exceeds the effect of a within-fibre swap of $P$ by at
least $\theta_{\mathrm{dir}}=0.05$. Sufficiency
($\Hent(C\mid\pa(C))\le\varepsilon_D\Hent(C)$), $|\pa(C)|\le k_{\max}$ and
resample stability remain required. Observational minimal sufficient sets
are still enumerated and reported as alternatives (Remark 2.7).
\item \textbf{Edge activity in the active derivation.} $D_b(x)$ keeps the
active nodes of Definition~\ref{def:active} and adds \emph{active edges}: $P\to C$ is
active on $x$ if the direct-effect swap of $P$ changes $C$ (or $r_b$) on at
least $\alpha$ times the fraction of the same partners whose own $C$ value
differs from $C(x)$. A node lies on the derivation if it is on an
active-edge path from an input to the output. Generality (Definition~\ref{def:generality})
counts inputs on whose derivation the node lies in this sense.
\item \textbf{Leaf rule.} In $\lambda$ and $\mu$ a node contributes $g(C)$ on
$x$ only if $C(x)\in\mathrm{Rg}(C)$ (Definition~\ref{def:range}); otherwise it is a
leaf on $x$ and contributes $0$. This is what makes ``each input is
effectively its own leaf'' (master prompt \S2.3.4) measurable: an exception
node whose value on $x$ is shared by fewer than $n_{\min}$ training inputs
does not ground $x$.
\item \textbf{Canonical realisers.} (i) In residual architectures the probe
realiser of a variable at site $s_\ell$ is decoded from what the block wrote,
$h_\ell-h_{\ell-1}$ (ridge decoder column space; $h_{\ell-1}$ is the same token
position at the previous exposed layer); at $s_0$ from the state. The ridge decoder is fitted inside the site's top principal subspace (at most
64 components; amendment 1, \texttt{prereg/PHASE\_1\_amend\_1.md}): with fewer
probe inputs than dimensions an unrestricted decoder interpolates along
near-null-variance directions that carry no causal effect. A probe
realiser is proposed only if its value map decodes its label with 2-fold
held-out accuracy $\ge0.9$ (proposal confers nothing; the filter only saves
interventions).
(ii) A candidate whose realiser lies inside a larger candidate (or input
realiser) at the same site whose partition refines it is a
\emph{coarsening} and is not a separate node. (iii) A candidate covered by the
span of two or more smaller, pairwise \emph{separable} (state-level modular)
quotient-passing candidates is a \emph{composite} and is not a node.
(iv) Equivalent candidates are merged; interventions use the realiser at the
node's earliest surviving site with the lowest quotient defect; members at
other sites are bookkeeping.
\item \textbf{The discovery fixed point.} Derivation failures are resolved
first, earliest site first; a node that fails derivation is suspended (it is
re-tested on every pass and reactivated if it passes once the pool changes).
Causal-use failures are resolved latest site first. Until the graph is
stable, causal use is judged on effective pairs (those for which $F_{\G}$
predicts a change) and a node with too few effective pairs is pending; the
final pass applies the full test (all-pairs IIA $\ge\tau_{\IIA}$, effective
IIA $\ge\tau_{\IIA}$, margin $\Delta_{\IIA}$ over norm-matched random realisers).
\item \textbf{Nulls and calibrations.} Random realisers are norm-matched
(each patch delta rescaled to the norm of the candidate's own delta).
Condition 1(i) is per behaviour: the realiser must matter more than the 95th
percentile of norm-matched random realisers for at least one behaviour.
$\kappa_{\max}$ is calibrated from the proposed $0.25$ to $0.6$ (the proposed
value rejects every non-trivial node of a two-input system); $\Hent(X)$ is
the sum of the input nodes' entropies.
\item \textbf{Output realisers.} A candidate whose partition is equivalent to a
behaviour's output partition is a realiser of that output node, not a
derivation step, exactly as input-equivalent candidates are realisers of
input nodes. Without this, the representation of the answer itself lies on
every derivation with generality equal to its behaviour's mass, and every
input of a single-behaviour suite is ``grounded'' through its own answer
(found on V2). The output node's parents are then the parents of that
representation.
\item \textbf{Interventional identification of derivation maps
(Proposition~\ref{prop:unseen}(b)).} When a child's parents include a node
and a function of it (positivity fails, Remark 2.7), every interchange on the
node produces a parent tuple never seen observationally, and $F_{\G}$ is
undefined there. Derivation maps $\psi_C$ are therefore fitted on the
observational tuples together with the tuples produced by the direct-effect
interventions of M1 (the patched parent at its source value, the others
clean or held). IIA is evaluated on fresh pairs; agreement at a tuple reached
both ways is the modularity condition of Proposition~\ref{prop:unseen}(b).
\item \textbf{Step-level leaf rule (amendment 2).} On $x$, a node contributes
$g(C)$ to $\lambda(x)$ and $\mu(x)$ only if $C(x)\in\mathrm{Rg}(C)$ and its parent
tuple $\pa(C)(x)\in\mathrm{Rg}(\pa(C))$ (Definition~\ref{def:range}); a value reached
through a parent tuple supported by fewer than $n_{\min}$ training inputs is a
lookup step. With the value-level rule alone, accurate realisers make an
exception node's stored labels ground memorized inputs whenever labels are
shared (the ``answer funnel'', P6b and amendment-1 P2). \emph{Consequence:} if
every node is computed from the full input identity, $\lambda=\mu=0$ everywhere
and structural memorization is undefined; reuse-based memorization presupposes
shared sub-structure in the inputs.
\end{enumerate}

\begin{remark}[What the amendments change in the statuses]
Propositions~\ref{prop:lattice}--\ref{prop:fca} concern the exact
observational objects and are unchanged. Proposition~\ref{prop:bottleneck} applies to
the active-edge DAG of (M2) with the node weights of (M3) verbatim.
\end{remark}

%============================================================================
\section{Propositions 2.7.1--2.7.3 of the master prompt}
%============================================================================

\begin{proposition}[CCGP and depth one; proved, \texttt{proofs/P7\_ccgp.tex}]
\label{prop:ccgp}
Bernardi et al.'s CCGP is not the depth-one case of depth: there are systems
with perfect depth-one grounding and CCGP $\le1/2$, and systems with CCGP $=1$
for a causally inert variable. With the decoder fixed to the system's own
readout and causal use added, readout-aligned CCGP $=1$ iff the depth-one
candidate passes the exact quotient and causal-use tests.
\end{proposition}

\begin{proposition}[Grounded behaviours generalize to lower inputs; proved
under (A1)--(A3), \texttt{proofs/P8\_grounded\_lower.tex}]\label{prop:lower}
If a set of principles is a vertex cut of the derivation on a new input,
the nodes below the cut are value-local, and the cut's joint value was taken
by a training input with the same derivation below the cut, the system's
output on the new input equals its output on that training input; it is
correct if the target factors through the principle values.
\end{proposition}

\begin{proposition}[Memorization and sideways generalization; proved,
\texttt{proofs/P9\_memorized\_sideways.tex}]\label{prop:memoside}
(a) A behaviour memorized at $\tau$ can generalize sideways (a private
compositional node used by one behaviour of a large suite). (b) If the
behaviour is leaf-memorized (every node active on a training input has an
empty training range), every new input is diagonal.
\end{proposition}

%============================================================================
\section{Summary of statuses}
%============================================================================

\begin{longtable}{p{0.33\textwidth}p{0.17\textwidth}p{0.42\textwidth}}
\toprule
Item & Status & Test \\
\midrule
\endhead
Prop.~\ref{prop:lattice} derivation as refinement & proved, unchecked & V1 \\
Prop.~\ref{prop:unique} parent uniqueness & proved, unchecked & V1, V3 \\
Prop.~\ref{prop:unseen} unseen combinations & proved, unchecked & K3 \\
Prop.~\ref{prop:reparam} reparametrisation & proved, unchecked & unit/property tests \\
Prop.~\ref{prop:welldef} well-definedness (exact) & proved, unchecked & V1--V4 \\
Prop.~\ref{prop:fca} FCA relation & proved, unchecked & V1, V2 \\
Prop.~\ref{prop:bottleneck} bottleneck & proved, unchecked & unit tests \\
Prop.~\ref{prop:order} order of scores & proved, unchecked & unit tests \\
Conj.~\ref{conj:stable} stability & conjecture & V1--V4, K1 \\
Conj.~\ref{conj:fca} approximate FCA & conjecture & V1, V2 \\
Conj.~\ref{conj:gamma} structural vs interventional & conjecture & V1, K1 \\
Conj.~\ref{conj:k2} depth predicts error (kill) & conjecture & K2 \\
Conj.~\ref{conj:memo} memorization verdict & conjecture & V2, K4, K9 \\
Conj.~\ref{conj:lower}--\ref{conj:diag} three moves & conjecture & V3, K3, K13 \\
Conj.~\ref{conj:creation} creation & conjecture & V5, K13 \\
Prop.~\ref{prop:ccgp} CCGP vs depth one & proved, unchecked & K2 (disagreement analysis) \\
Prop.~\ref{prop:lower} grounded $\Rightarrow$ lower correct & proved under (A1)--(A3), unchecked & K3, K6 \\
Prop.~\ref{prop:memoside} memorized vs sideways & proved (counterexample + leaf form), unchecked & V2, K3 \\
\bottomrule
\end{longtable}

\section*{Proposed thresholds (to be fixed in pre-registration)}

\begin{tabular}{lll}
\toprule
Symbol & Proposed value & Role \\
\midrule
$\varepsilon_Q$ & 0.10 & quotient defect \\
$\kappa_{\max}$, $h_{\min}$ & 0.25 (calibrated to 0.6, M6), 0.5 bit & compression, non-triviality \\
$\theta_{\mathrm{dir}}$ & 0.05 & direct-effect parent (M1) \\
$\varepsilon_D$, $k_{\max}$ & 0.05, 8 & parent sufficiency, parent count \\
$J_{\min}$, $\tau_E$ & 0.8, 0.9 & parent stability, edge interchange \\
$\tau_{\IIA}$, $\Delta_{\IIA}$ & 0.90, 0.20 & causal use, margin over null \\
$\varepsilon_{\mathrm{merge}}$ & 0.05 & equivalence \\
$\alpha$ & 0.5 & per-input activity \\
$\tau$ & 0.25 and 0.5 & principle levels \\
$\varepsilon_M$ & 0.10 & memorization verdict \\
$n_{\min}$, $k$ & 5, 10 & training range \\
$\varepsilon_{\mathrm{cov}}$ & 0.2 & diagonal coverage \\
\bottomrule
\end{tabular}

\appendix
\section{Proofs}
The proofs below are complete as written but have not been checked line by
line by a second reader; until that check is recorded they support
propositions, not theorems.

\subsection{Proof of Proposition~\ref{prop:lattice} (derivation as refinement)}
\label{proof:P1}
\textbf{Setting.} $X$ is finite and every $x\in X$ has $P(x)>0$ (restrict $X$
to the support of $P$; this is what ``modulo $P$-null sets'' means for
finite $X$). Concepts $C_1,\dots,C_k,C$ are maps from $X$ into finite sets.

\textbf{Claim.} There is a map $\psi$ with $C(x)=\psi(C_1(x),\dots,C_k(x))$
for all $x$ if and only if $\bigwedge_i\Part_{C_i}\le\Part_C$.

\textbf{Proof.} The blocks of the meet $\bigwedge_i\Part_{C_i}$ are the
non-empty sets $\{x:C_i(x)=t_i\ \forall i\}$, one for each tuple
$t=(t_1,\dots,t_k)$ in the image $T$ of $x\mapsto(C_1(x),\dots,C_k(x))$.

($\Rightarrow$) Let $x,x'$ lie in one block of the meet. Then
$C_i(x)=C_i(x')$ for every $i$, hence
$C(x)=\psi(C_1(x),\dots)=\psi(C_1(x'),\dots)=C(x')$, so $x,x'$ lie in one block
of $\Part_C$. Every block of the meet therefore lies inside a block of
$\Part_C$.

($\Leftarrow$) For $t\in T$ choose any $x_t$ with tuple $t$ and set
$\psi(t)=C(x_t)$; for $t\notin T$ set $\psi(t)$ arbitrarily. The choice does
not matter: all inputs with tuple $t$ form one block of the meet, which lies
in one block of $\Part_C$, so they share one value of $C$. Then for every $x$
with tuple $t$, $\psi(t)=C(x_t)=C(x)$. \hfill$\square$

\textbf{Corollary.} An exact parent set (an inclusion-minimal $S$ with
$\Hent(C\mid S)=0$) is exactly an inclusion-minimal $S\subseteq\mathcal{K}_{<s}$
with $\bigwedge_{C'\in S}\Part_{C'}\le\Part_C$, because for finite $X$ with
full support $\Hent(C\mid S)=0$ holds if and only if $C$ is a function of $S$
on $X$.

\subsection{Proof of Proposition~\ref{prop:unique} (uniqueness of the minimal parent
set) and of Proposition~\ref{prop:welldef}}
\label{proof:P2}
\textbf{Setting.} All variables are finite-valued functions of $X\sim P$.
$U=\{U_1,\dots,U_n\}$ is the candidate parent set; for $A\subseteq U$ write
$A$ also for the tuple of its members. \emph{Positivity}: for every
combination $u=(u_1,\dots,u_n)$ in $\prod_j\mathcal{V}_j$, where
$\mathcal{V}_j$ is the set of values $U_j$ takes with positive probability,
$P(U=u)>0$. Positivity of $U$ implies positivity of every sub-collection,
since a marginal of a positive law on a product set is positive on the
product of the marginal value sets. A set $S\subseteq U$ is a \emph{blanket}
for $C$ if $C\perp U\setminus S\mid S$.

\textbf{Lemma (intersection property with a deterministic target).} Let
$Y,W,Z$ be disjoint sub-collections of $U$ and assume positivity. If
$C\perp Y\mid(Z,W)$ and $C\perp W\mid(Z,Y)$, then $C\perp(Y,W)\mid Z$. No
positivity is required of $C$.

\textbf{Proof.} Fix values $z$, $y$, $w$. By positivity $P(y,w,z)>0$ for every
$y$ and $w$ in their value sets, so all conditionals below are defined. The
first hypothesis gives $P(c\mid y,z,w)=P(c\mid z,w)$ and the second gives
$P(c\mid y,z,w)=P(c\mid z,y)$. Hence $P(c\mid z,w)=P(c\mid z,y)$ for all
$y,w$. The left side does not depend on $y$ and the right side does not
depend on $w$, so both equal a function $q(c,z)$ only. Then
$P(c\mid z)=\sum_{y,w}P(c\mid y,z,w)P(y,w\mid z)=q(c,z)\sum_{y,w}P(y,w\mid z)
=q(c,z)$, and therefore $P(c\mid y,z,w)=q(c,z)=P(c\mid z)$ for all $y,w$, which
is $C\perp(Y,W)\mid Z$. \hfill$\square$

\textbf{Proof of Proposition~\ref{prop:unique}.} Let $S_1,S_2$ be blankets. Put
$Z=S_1\cap S_2$, $A=S_1\setminus S_2$, $B=S_2\setminus S_1$ and
$R=U\setminus(S_1\cup S_2)$; these are disjoint and cover $U$.
Since $S_1=Z\cup A$ is a blanket, $C\perp(B,R)\mid(Z,A)$.
Since $S_2=Z\cup B$ is a blanket, $C\perp(A,R)\mid(Z,B)$, and by weak union
$C\perp A\mid(Z,B,R)$. Apply the Lemma with $Y=A$ and $W=(B,R)$: the two
hypotheses are $C\perp A\mid(Z,B,R)$ and $C\perp(B,R)\mid(Z,A)$, both
established. Hence $C\perp(A,B,R)\mid Z$, i.e.\ $Z=S_1\cap S_2$ is a blanket.
Now let $S_1,S_2$ be inclusion-minimal blankets. Then $S_1\cap S_2$ is a
blanket contained in both, so by minimality $S_1=S_1\cap S_2=S_2$.
\hfill$\square$

\textbf{Link to Definition~\ref{def:parents}.} If $C$ is a function of $U$ (the exact
case of Condition 2, where some subset of $U$ is $0$-sufficient), then for
any $S\subseteq U$, $S$ is a blanket iff $\Hent(C\mid S)=0$. Indeed, if
$\Hent(C\mid S)=0$ then $C$ is a function of $S$ and is independent of
everything given $S$; conversely, if $C\perp U\setminus S\mid S$ then
$\Hent(C\mid S)=\Hent(C\mid S,U\setminus S)=\Hent(C\mid U)=0$. So under
positivity the exact parent set of Definition~\ref{def:parents} is unique.

\textbf{Proof of Proposition~\ref{prop:welldef}.} Nodes are equivalence classes of
accepted candidates, i.e.\ distinct partitions of $X$; since the accepted set
is fixed (assumption (c)), the node set is determined. Each class is assigned
the earliest site of its members under the fixed linear extension, which is
determined. For each node $C$ the candidate parents form a positive
collection by (b), so its exact parent set is unique by the result just
proved; hence the edge set is determined. Every edge goes from a node whose
site is strictly earlier to a node whose site is later in a fixed linear
order, so a directed cycle would require a site strictly earlier than
itself; the graph is acyclic. Finally, by Proposition~\ref{prop:reparam}, replacing a
value map by an invertible reparametrisation leaves the partition, hence the
node, and all conditional-independence statements, hence the edges,
unchanged. \hfill$\square$

\textbf{Caveat recorded with the proof.} Positivity fails whenever the
candidate set contains two members one of which is a function of the other
(for example a node together with its own parents), and equivalent
candidates violate it trivially; merging (Definition~\ref{def:equiv}) is therefore a
precondition, not a convenience. When positivity fails, distinct minimal
parent sets can exist (for $C=C_3=C_1\oplus C_2$ both $\{C_3\}$ and
$\{C_1,C_2\}$ are minimal). The programme then reports the set of all
minimal parent sets and every downstream quantity as a range over the
resulting graphs; the choice of a canonical candidate set (for example, the
latest-site antichain that is sufficient) is an open design decision for
Phase 1 and is tested on planted system V1.

\subsection{Proof of Proposition~\ref{prop:unseen} (unseen combinations)}
\label{proof:P3}
\textbf{(a)} Let $T=(C_1,\dots,C_k)$ be the parent tuple and $t$ a value with
$P(T=t)=0$, and let $\psi$ satisfy $C=\psi(T)$ $P$-a.s. Assume $|V|\ge2$
(otherwise $C$ is constant and fails $h_{\min}$). Choose $v\ne\psi(t)$ in $V$
and define $\psi'(t)=v$ and $\psi'=\psi$ elsewhere. Then
$P(\psi'(T)\ne\psi(T))\le P(T=t)=0$, so $C=\psi'(T)$ $P$-a.s., while
$\psi'(t)\ne\psi(t)$. No observational quantity distinguishes $\psi$ from
$\psi'$. \hfill$\square$

\textbf{(b)} This part is definitional and is stated so that its assumption
is explicit. \emph{Modularity} means: there is a map
$\Psi_{\mathrm{sys}}$, defined on the parent-value tuples reachable by joint
realiser patches, such that in every run (natural or patched) the value of
$C$ computed by the system equals $\Psi_{\mathrm{sys}}$ of the parent values
present in that run. Under modularity, a joint patch that produces parent
values $t$ yields a run in which $C$ takes the value $\Psi_{\mathrm{sys}}(t)$,
which is read off by $\varphi\circ R$ at the site of $C$. Hence
$\Psi_{\mathrm{sys}}(t)$ is observed, and it agrees with $\psi$ wherever
$P(T=\cdot)>0$. \hfill$\square$

\textbf{Remark.} Modularity is an assumption about the system, not a
consequence of Conditions 1--3; it is the assumption under which the
interchange test is evidence about the unseen combination, and it is itself
checked by testing that the patch does not alter $C$'s other parents
(Definition~\ref{def:parents}, test (iii)).

\subsection{Proof of Proposition~\ref{prop:reparam} (reparametrisation)}
\label{proof:P4}
Let $u\colon V\to V'$ be a bijection and $\tilde C=u\circ C$, with the same
site and realiser. The statement concerns population quantities, not the
behaviour of a particular finite-sample estimator.
\begin{enumerate}[label=(\arabic*)]
\item \emph{Fibre events.} $\tilde C(x)=\tilde C(x')$ iff $C(x)=C(x')$
because $u$ is injective, so $\Part_{\tilde C}=\Part_C$. The realiser patch
$f_{R\leftarrow x'}(x)$ does not involve $\varphi$. Hence $A_b^{=}$, $A_b$
and $\delta_Q$ are unchanged.
\item \emph{Entropies.} $\Hent(u(C))=\Hent(C)$ and
$\Hent(u(C)\mid S)=\Hent(C\mid S)$ for bijective $u$; likewise
$\Hent(C\mid S)$ is unchanged when a member of $S$ is reparametrised, since
conditioning depends only on the generated partition. So $\kappa$,
$h_{\min}$, sufficiency and minimality, and hence $\pa(\cdot)$, are unchanged.
\item \emph{Interchange.} With $\psi_{\tilde C}=u\circ\psi_C$ and each child's
map $\tilde\psi_{C''}$ obtained from $\psi_{C''}$ by precomposing $u^{-1}$ in
the $C$ argument, the high-level model satisfies
$\tilde F_{\G}(x;\mathrm{do}(\tilde C=u(v)))=F_{\G}(x;\mathrm{do}(C=v))$ for
every $v$, by direct substitution. The pairs with $\tilde C(x)\ne\tilde C(x')$
are the pairs with $C(x)\ne C(x')$, and the system-side patch is unchanged.
So $\IIA$ is unchanged.
\item \emph{Generality and depth.} $a_C(x)$ depends only on the fibre events
of (1) and on realiser patches, so $g$, $g_\#$, $\Pi_\tau$, $D(x)$,
$\lambda$, $\mu$ and $\gamma_\tau$ are unchanged.
\item \emph{Ranges.} For finite $V$, $\mathrm{Rg}(\tilde C)=u(\mathrm{Rg}(C))$
and $\tilde C(x^\ast)\in\mathrm{Rg}(\tilde C)$ iff $C(x^\ast)\in\mathrm{Rg}(C)$.
\end{enumerate}
\hfill$\square$

\textbf{Limit of the statement.} For a continuous value space, the
$\delta$-fibres of Definition~\ref{def:quotient} and the nearest-neighbour range of
Definition~\ref{def:range} are preserved by isometries of $(V,\rho)$ but not by general
homeomorphisms. The invariance claim for continuous $V$ is therefore made
only for isometries, and the choice of metric on $V$ is part of the
pre-registered specification of a concept.

\subsection{Proof of Proposition~\ref{prop:fca} (quotient condition and Formal Concept
Analysis)}
\label{proof:P5}
\textbf{Setting.} $\mathbb{K}_s=(X,M_s,I)$ with $X$ finite, $M_s$ a set of
maps $X\to\{0,1\}$ and $xIm\iff m(x)=1$. For $A\subseteq X$,
$A'=\{m:\forall x\in A,\ m(x)=1\}$; for $B\subseteq M_s$,
$B'=\{x:\forall m\in B,\ m(x)=1\}$. Extents are the sets $B'$ for
$B\subseteq M_s$ (equivalently, the sets $A$ with $A''=A$). $C$ is an exact
quotient for $M_s$: each $m$ is constant on each block of $\Part_C$.

\textbf{(a)} Let $E=B'$ be an extent and $x\in E$, and let $y\in[x]_C$. For
each $m\in B$, $m(y)=m(x)=1$, so $y\in B'=E$. Hence $E$ contains, with each of
its points, the whole $\Part_C$-block of that point, i.e.\ $E$ is a union of
blocks. \hfill$\square$

\textbf{(b)} $\{x\}''=\{x\}'{}'$ is the extent of the attribute set
$\{x\}'$; it contains $x$, so by (a) it contains $[x]_C$. \hfill$\square$

\textbf{(c)} Assume $M_s$ is closed under complement. Let $y\in\{x\}''$, so
$m(x)=1\Rightarrow m(y)=1$ for all $m\in M_s$. If $m(x)=0$ then
$(1-m)(x)=1$ with $1-m\in M_s$, so $(1-m)(y)=1$, i.e.\ $m(y)=0$. Hence
$m(y)=m(x)$ for every $m$, so $y$ lies in the kernel block $[x]_{M_s}$.
Conversely, if $y\in[x]_{M_s}$ then $y$ has every attribute $x$ has, so
$y\in\{x\}''$. Thus $\{x\}''=[x]_{M_s}$, and the object-concept extents are
exactly the blocks of $\Part_{M_s}$.
Since each $m$ is constant on $\Part_C$-blocks, two points in one
$\Part_C$-block have the same attribute vector, so $\Part_C\le\Part_{M_s}$:
every exact quotient refines $\Part_{M_s}$. The map
$x\mapsto(m(x))_{m\in M_s}$ is itself an exact quotient with partition
$\Part_{M_s}$, which is therefore the coarsest. Hence $C$ is the coarsest
exact quotient iff $\Part_C=\Part_{M_s}$, iff the blocks of $\Part_C$ are the
object-concept extents. \hfill$\square$

\textbf{Remark.} Without closure under complement, (c) fails: with a single
attribute $m$ and $m(x)=0$, $\{x\}'=\emptyset$ and $\{x\}''=X$, which is
strictly larger than the kernel block of $x$ whenever $m$ is not identically
zero. This is why the downstream context is built with nominal scaling.

\subsection{Proofs of Proposition~\ref{prop:bottleneck} and Proposition~\ref{prop:order}
(bottleneck characterisation and order of scores)}
\label{proof:P6}
\textbf{Conventions.} $D$ is a finite DAG with a set $I$ of input nodes, an
output node $o$, and weights $g\ge0$ on internal nodes (neither input nor
output). For a path $p$ from $I$ to $o$, $\max_{C\in p}g(C)$ is taken over
its internal nodes and is $0$ if it has none. $\mathcal{P}$ is the set of such
paths. If $\mathcal{P}=\emptyset$ we set $\lambda=0$ (an output disconnected
from the inputs is not grounded); $\mu$ is the maximum of $g$ over internal
nodes of $D$, and $0$ if there are none. A set $K$ of internal nodes is a
vertex cut if every $p\in\mathcal{P}$ contains a node of $K$. For $\tau>0$ let
$K_\tau=\{C:g(C)\ge\tau\}$.

\textbf{Proposition~\ref{prop:bottleneck}.} Let $\tau>0$ and $\mathcal{P}\ne\emptyset$.
$K_\tau$ is a vertex cut iff every $p\in\mathcal{P}$ contains an internal node
with $g\ge\tau$, iff $\max_{p}g\ge\tau$ for every $p$, iff
$\min_{p\in\mathcal{P}}\max_{C\in p}g(C)\ge\tau$ (the minimum over a finite
non-empty set is attained). Computation: process nodes in topological
order and set $M(v)=0$ for $v\in I$ and
$M(v)=\min_{u\to v}\max\bigl(M(u),\,g(u)\mathbf{1}[u\notin I]\bigr)$ otherwise.
By induction on the topological order, $M(v)$ is the minimum over paths from
$I$ to $v$ of the maximum weight of the internal nodes before $v$; so
$\lambda=M(o)$, computed with one pass over the edges. \hfill$\square$

\textbf{Proposition~\ref{prop:order}.} If $\mathcal{P}=\emptyset$ then $\lambda=0\le\mu$.
Otherwise, for any $p\in\mathcal{P}$, $\max_{C\in p}g(C)\le\mu$ because the
internal nodes of $p$ are internal nodes of $D$; hence
$\lambda=\min_p\max_{C\in p}g\le\mu$. If $\lambda(x)\ge\tau>0$ then
$\mu(x)\ge\tau$, so $x$ is not memorized at $\tau$.
For a behaviour: $\{\lambda\ge\tau\}\subseteq\{\mu\ge\tau\}$, so
$\Pr(\mu<\tau)\le1-\Lambda_\tau(b)$. If $b$ is grounded,
$\Lambda_\tau(b)\ge1-\varepsilon_M$, hence $\Pr(\mu<\tau)\le\varepsilon_M$,
which is $<1-\varepsilon_M$ whenever $\varepsilon_M<1/2$; so $b$ is not
memorized. Partly grounded is defined as neither, so the three verdicts are
mutually exclusive when $\varepsilon_M<1/2$. \hfill$\square$

\subsection{Proposition 2.7.1 of the master prompt: CCGP and depth one}
\label{proof:P7}
\textbf{Setting.} A single site with noiseless condition means: a finite set
of conditions $K$, a trial of condition $k$ has state $h=\mu_k\in\mathbb{R}^d$,
and conditions are equally weighted. A dichotomy $C\colon K\to\{0,1\}$. The
behaviour $b$ reads out $y(h)=\mathbf{1}[w^\top h>\theta]$ for a fixed
$w\ne0$, $\theta$. Bernardi et al.'s CCGP for $C$ with $m$ training conditions
per side is the mean, over all choices $S$ of $m$ conditions from each side,
of the accuracy on $K\setminus S$ of the max-margin linear classifier trained
on $\{\mu_k:k\in S\}$. The \emph{readout-aligned} CCGP replaces that
classifier by the system's own readout $y$. The depth-one graph is
$K\to C\to y$ with $C$ realised by $R=\mathrm{span}(w)$ and value map
$\varphi(z)=\mathbf{1}[z>\theta/\lVert w\rVert]$ on $z=\hat w^\top h$.

\textbf{Proposition P7.} (a) \emph{Depth-one grounding without CCGP.} There
is a system in which $y$ reports $C$ on every condition (so the candidate
passes the quotient test with $\delta_Q=0$, the causal-use test with
$\mathrm{IIA}=1$, and every input has $\lambda(x)=g(C)=1$), while CCGP with
$m=1$ is at most $1/2$.
(b) \emph{CCGP without a concept.} There is a system with CCGP $=1$ for a
dichotomy $C'$ whose realiser does not matter for any behaviour, so $C'$ fails
Condition 1(i) and contributes $0$ to every depth score.
(c) \emph{Aligned case.} Readout-aligned CCGP equals $1$ for every $m$ if and
only if $y(\mu_k)=C(k)$ for all $k$, which holds if and only if the depth-one
candidate passes the exact quotient and causal-use tests; then $\lambda(x)=g(C)=1$
on every input.
Hence Bernardi's CCGP is the depth-one case of the depth measure only after
the decoder is fixed to the system's own readout direction and causal use is
added; without these two changes neither score bounds the other.

\textbf{Proof.} (a) Take $d=2$, $C=1$ conditions $\mu_1=(1,0)$,
$\mu_2=(-3,10)$; $C=0$ conditions $\mu_3=(-1,0)$, $\mu_4=(3,-10)$; readout
$w=(1,\tfrac12)$, $\theta=0$. Then $w^\top\mu=1,2,-1,-2$, so $y=C$ on all four.
The realiser swap of $z=\hat w^\top h$ between two trials with the same $C$
leaves $y$ unchanged (the swapped coordinate stays on the same side of the
threshold), so $\delta_Q=0$; a swap between trials with different $C$ sets
$z$ to the partner's value and moves $y$ to the partner's $C$, which is what
the graph $C\to y$ predicts, so $\mathrm{IIA}=1$; $C$ is active on every
trial and lies on its only path, so $g(C)=1=\lambda(x)$. CCGP with $m=1$
averages four training pairs. The max-margin classifier for two points is the
perpendicular bisector. $\{\mu_1\}$ vs $\{\mu_3\}$: direction $(1,0)$; it
classifies $\mu_2$ ($-3$) as $0$ and $\mu_4$ ($3$) as $1$: accuracy $0$.
$\{\mu_2\}$ vs $\{\mu_4\}$: direction $(-6,20)$ through the origin; $\mu_1$
scores $-6$ (wrong), $\mu_3$ scores $6$ (wrong): accuracy $0$.
$\{\mu_1\}$ vs $\{\mu_4\}$: direction $\mu_1-\mu_4=(-2,10)$, midpoint
$(2,-5)$, offset $(-2,10)\cdot(2,-5)=-54$; $\mu_2$ scores $106+54>0$ (right),
$\mu_3$ scores $2+54>0$ (wrong): accuracy $1/2$. $\{\mu_2\}$ vs $\{\mu_3\}$
is symmetric: accuracy $1/2$. So CCGP $=(0+0+\tfrac12+\tfrac12)/4=1/4\le1/2$.
(b) Take $d=2$, $\mu_k=(C'(k),\,\epsilon_k)$ with arbitrary small $\epsilon_k$
and a behaviour that reads only the second coordinate. Every training split
yields a classifier whose direction is dominated by the first coordinate, and
all held-out conditions are classified correctly (both sides are separated
by $x_1=\tfrac12$ with margin $\tfrac12$ against $|\epsilon|$ small), so
CCGP $=1$. Swapping the first coordinate never changes any readout, so
$A_b(C')=0$ for every behaviour and Condition 1(i) fails.
(c) With the decoder fixed to $y$, the accuracy on held-out conditions is
the fraction of $k\in K\setminus S$ with $y(\mu_k)=C(k)$. It is $1$ for every
$S$ iff $y(\mu_k)=C(k)$ for all $k$ (every condition is held out by some
$S$). If $y(\mu_k)=C(k)$ for all $k$, the swap argument of (a) gives
$\delta_Q=0$ and $\mathrm{IIA}=1$, and $\lambda=g=1$. Conversely if
$y(\mu_k)\ne C(k)$ for some $k$, a swap from a trial of condition $k$ into a
trial of the other side moves $y$ to $y(\mu_k)$, not to $C(k)$, so
$\mathrm{IIA}<1$. \hfill$\square$

\textbf{Measurable prediction.} On real models, CCGP and depth disagree in
two identifiable ways: variables that are geometrically abstract but unused
(high CCGP, failing Condition 1 or 3) and variables that are used but not
parallel across conditions (low CCGP, high depth). The K2 contrast reports
the per-input agreement of the two scores and the share of disagreements of
each kind; that is what decides whether depth carries information CCGP lacks.

\subsection{Proposition 2.7.2 of the master prompt: grounded behaviours and lower inputs}
\label{proof:P8}
\textbf{Setting and assumptions.} Fix a behaviour $b$, its concept graph
$\G$ and a set of principles $\Pi\subseteq\Pi_\tau$. Let $T$ be the training
inputs and $x^\ast$ a new input. Write $\mathrm{Desc}(\Pi)$ for the nodes of
$D_b(x^\ast)$ that are strictly downstream of $\Pi$ on a path to the output
node $o$ (including $o$).
\begin{enumerate}[label=(A\arabic*)]
\item \emph{Grounding on $x^\ast$.} $\Pi$ is a vertex cut of $D_b(x^\ast)$
between the input nodes and $o$.
\item \emph{Value locality downstream of the cut.} For every
$C\in\mathrm{Desc}(\Pi)$ and every $x\in T\cup\{x^\ast\}$,
$C(x)=\psi_C(\pa_{x}(C)(x))$, where $\pa_x(C)$ are the parents of $C$ that are
active on $x$, and $\psi_C$ is a fixed function (the system's derivation map,
which the exact derivation and interchange conditions identify on the
parent tuples seen in training).
\item \emph{$x^\ast$ is lower for $\Pi$.} There is a training input
$x_t\in T$ with $D_b(x_t)=D_b(x^\ast)$ on $\mathrm{Desc}(\Pi)\cup\Pi$ and
$\Pi(x_t)=\Pi(x^\ast)$ (the joint tuple of principle values is in range).
\end{enumerate}

\textbf{Proposition P8.} Under (A1)--(A3), $f_b(x^\ast)=f_b(x_t)$. Consequently, if
training error is zero on $x_t$ and the task's target factors through the
principle values ($\mathrm{target}(x)=\rho(\Pi(x))$ for some $\rho$ on
$T\cup\{x^\ast\}$), the system answers $x^\ast$ correctly.

\textbf{Proof.} Order $\mathrm{Desc}(\Pi)$ topologically. Claim: every
$C\in\mathrm{Desc}(\Pi)$ has $C(x^\ast)=C(x_t)$. Take $C$ and suppose the
claim holds for all earlier members. Each active parent $P$ of $C$ on $x^\ast$
lies on an active path to $o$ through $C$. If $P$ were neither in $\Pi$ nor in
$\mathrm{Desc}(\Pi)$, the active path from an input to $P$ (which exists
because $P\in D_b(x^\ast)$ lies on an input-to-output path) followed by
$P\to C\to\dots\to o$ would avoid $\Pi$: it cannot meet $\Pi$ after $C$
because $C$ is already downstream of $\Pi$ and the graph is acyclic, and it
does not meet $\Pi$ before $P$ by the choice of $P$. That contradicts (A1).
So every active parent of $C$ is in $\Pi$ or is an earlier member of
$\mathrm{Desc}(\Pi)$; by (A3) and the induction hypothesis its value on
$x^\ast$ equals its value on $x_t$, and by (A3) the active parent sets
agree. By (A2), $C(x^\ast)=\psi_C(\cdot)=C(x_t)$. Applying the claim to
$C=o$ gives $f_b(x^\ast)=f_b(x_t)$. If $f_b(x_t)$ is correct and the target
factors through $\Pi$, then $\mathrm{target}(x^\ast)=\rho(\Pi(x^\ast))=
\rho(\Pi(x_t))=\mathrm{target}(x_t)=f_b(x_t)=f_b(x^\ast)$. \hfill$\square$

\textbf{What the assumptions cost.} (A2) is the substantive one: it says the
nodes below the principles are computed from their parents' values only,
also on $x^\ast$. It is measurable on $x^\ast$ by within-fibre swaps and by
edge interventions from training partners (the interchange identification of
Proposition~\ref{prop:unseen}(b)). (A3) is the definition of a lower input,
restricted to the cut and below it; nodes upstream of the cut may be out of
range. The factorisation of the target is a property of the task and the
principles, not of the system. It fails exactly when a principle is a
shortcut (a high-generality node that predicts the target on the training
distribution only), and the proposition then predicts confident errors on
lower inputs, the K6 shortcut signature.

\textbf{Measurable prediction.} Among held-out inputs that are lower for the
principle cut, errors occur only where (A2) fails on the input (detected by
interventions on that input) or where the target does not factor through the
principles; errors on lower inputs with (A2) verified are therefore
concentrated in behaviours whose principles are known shortcuts. Tested in
K3 (moves) and K6 (shortcuts).

\subsection{Proposition 2.7.3 of the master prompt: memorization and sideways generalization}
\label{proof:P9}
The statement ``a memorized behaviour has no sideways generalization'' is
false for memorization as defined (Definition~\ref{def:memo}); it is true for a
stronger, leaf-level form.

\textbf{Proposition P9.} (a) \emph{Counterexample.} For every $\tau>0$ there
is a system and a behaviour $b$ that is memorized at $\tau$, has a non-empty
set of sideways inputs, and answers every sideways input correctly.
(b) \emph{Leaf-level memorization.} Say $b$ is \emph{leaf-memorized} if every
non-input node $C$ of $\G$ that is active on some training input of $b$ has an
empty training range, $\mathrm{Rg}(C)=\emptyset$ (no value of $C$ is taken on
$n_{\min}$ or more distinct training inputs). Then no new input is lower or
sideways for $b$: every new input whose derivation contains a non-input node
is diagonal, and so is every new input whose output is not mediated by the
graph.

\textbf{Proof.} (a) Let $m>1/\tau$ behaviours $b_1,\dots,b_m$ with equal weight
$1/m<\tau$. Behaviour $b_1$ computes $C=(A+B)\bmod q$ through a private node $C$
used by no other behaviour, from nodes $A,B$ computed from disjoint inputs,
and reads $y=C$; behaviours $b_2,\dots,b_m$ read other inputs through other
private nodes. Training withholds a set $W$ of $(A,B)$ combinations while
every value of $A$ and of $B$ is seen ($q\ge3$, $|W|$ small). The block that
computes $C$ is the full addition table (the compositional variant of planted
system V3). Every node on the derivation of $b_1$ ($A$, $B$, $C$) is used only
by $b_1$, so $g\le 1/m<\tau$, $\mu(x)<\tau$ on every input, and $b_1$ is
memorized at $\tau$. An input with $(A,B)\in W$ has all node values in range
and the parent tuple of $C$ unseen, so it is sideways; the system writes
$(A+B)\bmod q$ there, which is correct.
(b) Let $x^\ast$ be new. If $D_b(x^\ast)$ has a non-input node $C$, then
$C(x^\ast)$ is a value of $C$ and $\mathrm{Rg}(C)=\emptyset$, so
$C(x^\ast)\notin\mathrm{Rg}(C)$; by Definition~\ref{def:moves} $x^\ast$ is
neither lower nor sideways, hence diagonal (operationally, the out-of-range
value is reported). If $D_b(x^\ast)$ has no non-input node, the output is not
mediated by the graph, which is diagonal by the operational rule of
Definition~\ref{def:moves}(iii) (graph mediation share $0<1-\varepsilon_{\mathrm{cov}}$).
\hfill$\square$

\textbf{Consequences.} The structural memorization verdict is a statement
about reuse across the behaviour suite, not about compositionality inside one
behaviour: (a) shows it can label a compositional but unshared computation as
memorized. That is a design property of the definition, and it fixes what the
K4/K9 memorization claims may say: a model judged memorized is predicted to
have no \emph{shared} principles, not to fail on every sideways input. The
leaf-level form (b) is the one that predicts failure, through
Conjecture~\ref{conj:diag}.

\textbf{Measurable prediction.} (i) On systems trained with a planted
memorized subset (V2 and its trained variant), inputs whose derivation is
leaf-memorized are classified diagonal by $r.\mathrm{move}$ when perturbed to
unseen neighbours, and fail. (ii) The counterexample family exists in real
suites: behaviours judged memorized at $\tau$ whose private nodes pass the
interchange test on unseen tuples generalize sideways; K3 reports how often.


\section*{References}
Full, verified citations are in \texttt{docs/novelty.md}. Those used here:
R.~Wille (1982), Restructuring lattice theory; B.~Ganter and R.~Wille (1999),
Formal Concept Analysis; A.~Geiger, H.~Lu, T.~Icard and C.~Potts (2021),
Causal abstractions of neural networks; A.~Geiger et al.\ (2025), Causal
abstraction: a theoretical foundation for mechanistic interpretability;
K.~Chalupka, P.~Perona and F.~Eberhardt (2015), Visual causal feature
learning; S.~Bernardi et al.\ (2020), The geometry of abstraction in the
hippocampus and prefrontal cortex; J.~Pearl (1988), Probabilistic Reasoning
in Intelligent Systems (Markov boundaries and the intersection property).

\end{document}
